\documentclass{article}
\usepackage{amsmath, amssymb}
\usepackage{bm} % bold math
\usepackage{enumitem}
\usepackage{svg}
\usepackage[a4paper, margin=0.5in]{geometry}
% Include hyperref as the last package (except for a few exceptions)
\usepackage[colorlinks=true,
linkcolor=blue,
citecolor=blue,
urlcolor=blue,
pdfpagemode=FullScreen,
filecolor=blue,]{hyperref}


\begin{document}
\begin{titlepage}
\vspace*{0.5cm}
\begin{center}
\includesvg{LOGO}
\vspace*{2cm}
    \section*{Derivation Using the Explicit Solution of Fick's Second Law}
    \vfill
    November 2023
\end{center}
\end{titlepage}
\newpage

\textbf{1. Fick's Second Law:}

Starting with Fick's Second Law \cite{Fick1855}:

\begin{equation}\label{eq:con:1}
\frac{\partial C(x,t)}{\partial t} = D\,\frac{\partial^2 C(x,t)}{\partial x^2}\,
\end{equation}

The boundary conditions are as follows:
\begin{itemize}[leftmargin=2.5em]
    \item \textbf{Initial condition:} \( C(x,0) = C_0 \) for \( x \ge 0 \).
    \item \textbf{Far-field condition:} \( C(x \to \infty, t) = C_0 \).
    \item \textbf{Boundary condition at \( x=0 \):} A constant flux is imposed:
    
    \begin{equation}\label{eq:con:2}
    -D\,\left.\frac{\partial C}{\partial x}\right|_{x=0} = \frac{I}{F\,A}\,
    \end{equation}
    
    where \(I\) is the applied current, \(F\) is Faraday's constant, and \(A\) is the electrode area.
\end{itemize}

\vspace{1ex}
\textbf{2. Introduce the Similarity Variable:}

Let

\begin{equation}\label{eq:con:3}
\eta = \frac{x}{2\sqrt{Dt}}\,
\end{equation}

This substitution transforms the diffusion equation into an ordinary differential equation in terms of \(\eta\).

\vspace{1ex}
\textbf{3. Solve the ODE:}

The solution that satisfies the above conditions is

\begin{equation}\label{eq:con:4}
C(x,t) = C_0 - \frac{I}{F\,A\,D}\Biggl[\,x\,\operatorname{erfc}\Bigl(\frac{x}{2\sqrt{Dt}}\Bigr) - \frac{2\sqrt{Dt}}{\sqrt{\pi}}\,\exp\Bigl(-\frac{x^2}{4Dt}\Bigr)\Biggr]\,
\end{equation}

\vspace{1ex}
\textbf{4. Evaluation at the Electrode Surface (\( x=0 \)):}

Since \(\operatorname{erfc}(0)=1\) and the term \(x\,\operatorname{erfc}(\cdot)\) vanishes when \( x=0 \), we have

\begin{equation}\label{eq:con:5}
C(0,t) = C_0 - \frac{2I}{F\,A\,D}\,\sqrt{\frac{Dt}{\pi}}\,
\end{equation}

The magnitude of the change is:

\begin{equation}\label{eq:con:6}
\Delta C = \Bigl|C(0,t)-C_0\Bigr| = \frac{2I}{F\,A}\,\sqrt{\frac{t}{\pi\,D}}\,
\end{equation}

For a pulse of duration \(\tau\) (i.e. \(t=\tau\)):

\begin{equation}\label{eq:con:7}
\Delta C = \frac{2I}{F\,A}\,\sqrt{\frac{\tau}{\pi\,D}}\,
\end{equation}

\vspace{1ex}
\textbf{5. Link with the Voltage Response:}

Assume a linear relationship between the electrode potential and the concentration change,

\begin{equation}\label{eq:con:8}
\Delta E_\tau = \Bigl(\frac{dE}{dC}\Bigr)\Delta C\,
\end{equation}

Thus,

\begin{equation}\label{eq:con:9}
\Delta E_\tau = \Bigl(\frac{dE}{dC}\Bigr)\frac{2I}{F\,A}\,\sqrt{\frac{\tau}{\pi\,D}}\,
\end{equation}

Rearrange to solve for \(D\). Isolating the square-root gives:

\begin{equation}\label{eq:con:10}
\sqrt{\frac{1}{D}} = \frac{\Delta E_\tau}{\frac{dE}{dC}}\,\frac{F\,A}{2I}\,\sqrt{\frac{\pi}{\tau}}\,
\end{equation}

Squaring both sides yields

\begin{equation}\label{eq:con:11}
\frac{1}{D} = \frac{(\Delta E_\tau)^2}{\Bigl(\frac{dE}{dC}\Bigr)^2}\Bigl(\frac{F\,A}{2I}\Bigr)^2\,\frac{\pi}{\tau}\,
\end{equation}

or equivalently,

\begin{equation}\label{eq:con:12}
D = \frac{4I^2\,\tau}{\pi\,F^2\,A^2\,(\Delta E_\tau)^2}\Bigl(\frac{dE}{dC}\Bigr)^2\,
\end{equation}

\vspace{1ex}
\textbf{6. Express \(\frac{dE}{dC}\) Using Equilibrium Measurements:}

In a GITT experiment, after a current pulse the electrode relaxes to a new equilibrium potential. The overall (steady-state) voltage change \(\Delta E_s\) is associated with the overall concentration change \(\Delta C_\text{tot}\) in the active material. Using Faraday’s law, we have

\begin{equation}\label{eq:con:13}
\Delta C_\text{tot} = \frac{I\,\tau}{F}\,\frac{M}{m\,V_M}\,
\end{equation}

and, assuming

\begin{equation}\label{eq:con:14}
\Delta E_s = \Bigl(\frac{dE}{dC}\Bigr)\Delta C_\text{tot}\,
\end{equation}

we obtain

\begin{equation}\label{eq:con:15}
\frac{dE}{dC} = \Delta E_s\,\frac{F\,m\,V_M}{I\,\tau\,M}\,
\end{equation}

\vspace{1ex}
\textbf{7. Final Expression for \(D\):}

Substitute the expression for \(\frac{dE}{dC}\) into the equation for \(D\):

\begin{equation}\label{eq:con:16}
D = \frac{4I^2\,\tau}{\pi\,F^2\,A^2\,(\Delta E_\tau)^2} \left[\Delta E_s\,\frac{F\,m\,V_M}{I\,\tau\,M}\right]^2\,
\end{equation}

Simplify step-by-step:
\begin{enumerate}[label=(\roman*)]
    \item Square the bracket:
    \begin{equation}\label{eq:con:17}
    \left[\Delta E_s\,\frac{F\,m\,V_M}{I\,\tau\,M}\right]^2 = (\Delta E_s)^2\,\frac{F^2\,m^2\,V_M^2}{I^2\,\tau^2\,M^2}\,
    \end{equation}
    \item Multiply by the prefactor:
    \begin{equation}\label{eq:con:18}
    D = \frac{4I^2\,\tau}{\pi\,F^2\,A^2\,(\Delta E_\tau)^2} \cdot \frac{(\Delta E_s)^2\,F^2\,m^2\,V_M^2}{I^2\,\tau^2\,M^2}\,
    \end{equation}
    \item Cancel the common factors:
    \begin{equation}\label{eq:con:19}
    D = \frac{4\,m^2\,V_M^2\,(\Delta E_s)^2}{\pi\,A^2\,M^2\,\tau\,(\Delta E_\tau)^2}\,
    \end{equation}
\end{enumerate}
Thus, the final GITT expression for the diffusion coefficient is:

\begin{equation}\label{eq:con:20}
\bm{D = \frac{4}{\pi\,\tau}\left(\frac{m\,V_M}{M\,A}\right)^2\left(\frac{\Delta E_s}{\Delta E_\tau}\right)^2}\,
\end{equation}
\newpage
\bibliographystyle{ieeetr}
\bibliography{references}

\end{document}
